\textbf{定义与公式：} 通径是过焦点且垂直于椭圆长轴、双曲线实轴或抛物线对称轴的弦段。椭圆a>b>0，双曲线a,b>0，通径整长 $L=2b^2/a$。抛物线 $y^2=\pm2px$ 或 $x^2=\pm2py$，p>0为焦准距，整长L=2p。反射、旋转、平移不改长度。
\textbf{最低门槛证明：} 横向模型代焦点横坐标：椭圆与双曲线均得y²=b⁴/a²，两纵坐标±b²/a，距离2b²/a；抛物线代x=p/2得y=±p，距离2p。
\textbf{几何证明：} 通径端点的准距是|c-a²/c|，乘焦准比c/a得到半通径b²/a。抛物线两平行线间距p，等距定义给半通径p；对称后乘2。
\textbf{边界：} 只过焦点而未垂直，不一定是通径；公式仍取决于尺寸，不是只由形状决定。例x²/25+y²/9=1通径18/5，但同焦点斜弦y=x-4长90/17。
